\exercisesheader{}

% 1

\eoce{\qt{Migrain dan akupunktur,
    Bagian I\label{migraine_and_acupuncture_intro}}
Migrain merupakan jenis sakit kepala yang sangat nyeri,
dan sebagian pasien berupaya mengatasinya dengan akupunktur.
Untuk menentukan apakah akupunktur meredakan nyeri
migrain, para peneliti melakukan uji terkontrol acak
yang melibatkan 89 perempuan yang didiagnosis menderita migrain.
Mereka dialokasikan secara acak ke salah satu dari dua kelompok:
perlakuan atau kontrol.
Sebanyak 43 pasien dalam kelompok perlakuan menerima akupunktur
yang dirancang secara khusus untuk menangani migrain.
Sebanyak 46 pasien dalam kelompok kontrol menerima akupunktur plasebo
(penusukan jarum pada lokasi yang bukan titik akupunktur).
24 jam setelah menerima akupunktur, pasien ditanya
apakah mereka bebas nyeri.
Hasilnya dirangkum dalam tabel kontingensi
di bawah ini.\footfullcite{Allais:2011}

\noindent\begin{minipage}[l]{0.4\textwidth}
\begin{tabular}{ll  cc c} 
			                         		&           & \multicolumn{2}{c}{\textit{Bebas nyeri}} \\
\cline{3-4}
			                        	 	&			& Ya 	& Tidak 	                  & Jumlah \\
\cline{2-5}
							& Perlakuan 	& 10	 	& 33		                  & 43 \\
\raisebox{1.5ex}[0pt]{\emph{Kelompok}} & Kontrol	 	& 2	 	& 44 	 	                  & 46 \\
\cline{2-5}
							& Jumlah		& 12		& 77		                  & 89
\end{tabular}
\end{minipage}

\begin{quote}
\small\textit{Deskripsi tekstual: Pada skema telinga dalam artikel asli,
titik M berada di dekat bagian bawah depan daun telinga, sedangkan titik S
berada di bagian tengah atas. Artikel tersebut membedakan area akupunktur
yang sesuai (M) dan yang tidak sesuai (S) untuk menangani serangan migrain.
Gambar pihak ketiga tidak disertakan dalam edisi turunan ini.}
\end{quote}
\begin{parts}
\item Berapa persen pasien dalam kelompok perlakuan yang bebas nyeri 24 jam
setelah menerima akupunktur?
\item Berapa persen pasien dalam kelompok kontrol yang bebas nyeri?
\item Kelompok mana yang memiliki persentase pasien bebas nyeri lebih tinggi 24 jam
setelah menerima akupunktur?
\item Temuan Anda sejauh ini mungkin menunjukkan bahwa akupunktur merupakan terapi
yang efektif bagi semua orang yang menderita migrain. Namun, ini bukan satu-satunya
kesimpulan yang dapat ditarik berdasarkan temuan Anda sejauh ini. Apa satu
penjelasan lain yang mungkin atas perbedaan teramati antara persentase
pasien yang bebas nyeri 24 jam setelah menerima akupunktur pada kedua kelompok?
\end{parts}
}{}

% 2

\eoce{\qt{Sinusitis dan antibiotik,
    Bagian I\label{sinusitis_and_antibiotics_intro}} 
Para peneliti yang membandingkan efek pengobatan antibiotik untuk sinusitis akut
dengan penanganan simtomatik mengalokasikan secara acak 166 orang dewasa yang didiagnosis
menderita sinusitis akut ke salah satu dari dua kelompok: perlakuan atau kontrol. Peserta
penelitian menerima amoksisilin (antibiotik) selama 10 hari
atau plasebo yang rupa dan rasanya serupa. Plasebo tersebut terdiri atas
penanganan simtomatik seperti asetaminofen, dekongestan hidung, dan sebagainya.
Pada akhir periode 10 hari, pasien ditanya apakah
gejala mereka membaik.
Distribusi jawaban dirangkum di bawah ini. 
\footfullcite{Garbutt:2012}
\begin{center}
\begin{tabular}{ll  cc c} 
                                    			&			& \multicolumn{2}{c}{\textit{Perbaikan gejala}} \\
                                    			&			& \multicolumn{2}{c}{\textit{menurut laporan pasien}} \\
\cline{3-4}
			                        		&			& Ya 	& Tidak 	& Jumlah \\
\cline{2-5}
							& Perlakuan 	& 66		& 19		& 85 \\
\raisebox{1.5ex}[0pt]{\emph{Kelompok}}	& Kontrol		& 65		& 16 		& 81 \\
\cline{2-5}
							& Jumlah		& 131	& 35		& 166
\end{tabular}
\end{center}
\begin{parts}
\item Berapa persen pasien dalam kelompok perlakuan yang mengalami perbaikan
gejala?
\item Berapa persen pasien dalam kelompok kontrol yang mengalami perbaikan
gejala?
\item Kelompok mana yang memiliki persentase pasien dengan perbaikan
gejala lebih tinggi?
\item
    Temuan Anda sejauh ini mungkin menunjukkan adanya perbedaan nyata
    antara efektivitas antibiotik dan plasebo
    dalam memperbaiki gejala sinusitis.
    Namun, ini bukan satu-satunya kesimpulan yang
    dapat ditarik berdasarkan temuan Anda sejauh ini.
    Apa satu penjelasan lain yang mungkin atas perbedaan teramati
    antara persentase pasien dalam kelompok perlakuan
    antibiotik dan plasebo yang mengalami
    perbaikan gejala sinusitis?
\end{parts}
}{}
